% !TeX root = ../../Book2.tex
% 纲要标题：### 17.3 Skolem–Noether 定理
% 纲要位置：计划纲要.md，第 1609 行。

\section{Skolem–Noether 定理}
\label{b2:sec:1703}

中心化子定理比较一个子代数与同它交换的部分。现在比较同一个简单代数嵌入 \(A\) 的两种方式：结论是，它们只相差 \(A\) 中一个可逆元的内共轭。证明不从方程组中猜这个可逆元，而是把两种嵌入看成两个模结构，再由简单 Artin 环的模分类构造所需同构。

% 纲要要点（供目录核对）：
%
% * 单子代数的嵌入及内自同构；两种嵌入只改变左作用，单 Artin 环的一种简单模类型使维数决定重数
% * 中心单代数的自同构；从右正则模同构在单位元的像构造共轭元，逐个自同构可提升不保证存在群截面
% * 适用假设与非单代数的反例；非单源的不同分量重数及非中心环境的外自同构分别说明失败原因
%

\subsection{单子代数的嵌入与内自同构}
\label{b2:subsec:1703:01}

\begin{theorem}[Skolem–Noether]\label{b2:thm:skolem-noether}
设 \(k\) 为域，\(A\) 为中心单 \(k\) 代数，\(B\) 为非零有限维简单含幺 \(k\) 代数。对任意两个保幺 \(k\) 代数同态 \(\varphi,\psi:B\to A\)，存在 \(u\in A^\times\)，使
\[
\psi(b)=u\cdot\varphi(b)u^{-1}\qquad(b\in B).
\]
这些同态因 \(B\) 简单而自动单射；\(Z(B)/k\) 不必可分。这个版本也见\cite[Main Theorem 7.7.1]{Voight2021QuaternionAlgebras}。
\end{theorem}
\begin{proof}
把 \(B\) 的左作用与 \(A\) 的右正则作用放在同一个代数 \(E=B\otimes_kA^{\mathrm{op}}\) 中，便能只改变嵌入给出的左作用而保留右作用。令 \(L=Z(B)\)。与\cref{b2:thm:csa-centralizer}的证明相同，
\[
E\cong B\otimes_L(L\otimes_kA^{\mathrm{op}})
\]
是两个中心单 \(L\) 代数的张量积，因此为单 Artin 环。这一步只用任意标量扩张的中心单性，不要求 \(L/k\) 可分。

在同一个 \(k\) 向量空间 \(A\) 上定义两种左 \(E\) 模结构：
\[
(b\otimes a^{\mathrm{op}})\cdot_\varphi x=\varphi(b)xa,
\qquad
(b\otimes a^{\mathrm{op}})\cdot_\psi x=\psi(b)xa.
\]
左右乘法交换，所以模公理成立。由\cref{b2:prop:simple-artinian}写 \(E\cong M_r(D)\)，再由\cref{b2:thm:artin-wedderburn}及\cref{b2:prop:semisimple-multiplicity}的矩阵块分类，简单左模只有 \(S=D^r\) 一种类型。两种模都是有限个 \(S\) 的直和，且底层 \(k\) 空间都是 \(A\)，因此
\[
A_\varphi\cong S^{\oplus m}\cong A_\psi,\qquad
m=\frac{\dim_kA}{\dim_kS}.
\]
这里可以只凭维数确定重数，正因为 \(E\) 是单 Artin 环，只有一种简单左模。一般半单环可能有多种简单类型，它们的重数不能只由总维数恢复。于是存在 \(E\) 模同构 \(F:A_\varphi\to A_\psi\)。

\(E\) 线性首先保证 \(F(xa)=F(x)a\)，即 \(F\) 是右正则模之间的同构。右正则模由 \(1\) 生成，因此令 \(u=F(1)\)，便有 \(F(x)=F(1\cdot x)=ux\)。逆同构也由一的像决定，故为左乘某个 \(v\in A\)；两次复合为恒等在 \(1\) 处给出 \(vu=uv=1\)，所以 \(u\) 可逆。再由 \(B\) 的作用相容性，
\[
u\cdot\varphi(b)=F\bigl(\varphi(b)\bigr)=\psi(b)F(1)=\psi(b)u.
\]
右乘 \(u^{-1}\) 得到所需公式。最后，同态的核是 \(B\) 的双边理想，且因保持单位元不能为 \(B\)，所以核为零。
\end{proof}

这称为\term[Skolem-Noether-dingli]{Skolem–Noether 定理}[Skolem--Noether theorem]。若 \(B_1,B_2\subseteq A\) 为简单 \(k\) 子代数，且给定同构 \(\alpha:B_1\to B_2\)，把 \(B_1\) 的包含与 \(\alpha\) 后接 \(B_2\) 的包含作为两种嵌入，就得到 \(\alpha\) 可延伸为 \(A\) 的一个内自同构。所选 \(u\) 一般不唯一；若 \(u,v\) 实现同一嵌入之间的共轭，则 \(v^{-1}u\) 恰好与 \(\varphi(B)\) 的全部元素交换。

\subsection{中心单代数的自同构}
\label{b2:subsec:1703:02}

\begin{corollary}\label{b2:cor:csa-automorphisms-inner}
设 \(k\) 为域，\(A\) 为中心单 \(k\) 代数，则 \(A\) 的每个 \(k\) 代数自同构都是内自同构，并有群同构
\[
\operatorname{Aut}_{k\text{-alg}}(A)\cong A^\times/k^\times.
\]
特别地，对每个整数 \(n\ge1\)，\(\operatorname{Aut}_{k\text{-alg}}\bigl(M_n(k)\bigr)\cong\operatorname{PGL}_n(k)\)，其中
\(\operatorname{PGL}_n(k)=\operatorname{GL}_n(k)/k^\times I_n\)。
\end{corollary}
\begin{proof}
在\cref{b2:thm:skolem-noether}中取 \(B=A\)，两个嵌入分别为恒等映射和给定自同构，得到满群同态
\(A^\times\to\operatorname{Aut}_{k\text{-alg}}(A)\)、\(u\mapsto(a\mapsto uau^{-1})\)。它的核由与全部 \(A\) 交换的可逆元组成，即 \(Z(A)^\times=k^\times\)。因此把商类 \(uk^\times\) 送到该共轭自同构，给出所述商群同构。矩阵代数情形直接代入。
\end{proof}
\symidx{\(\operatorname{PGL}_n(k)\)：一般线性群模去非零标量矩阵所得的商群}

例如，若 \(\operatorname{char}k\ne2\)，\(d\in k^\times\) 非平方，矩阵
\[
T=\begin{pmatrix}0&d\\1&0\end{pmatrix}
\]
给出二次域 \(k(\sqrt d)\) 的嵌入。其非平凡自同构 \(\sqrt d\mapsto-\sqrt d\) 由 \(P=\operatorname{diag}(1,-1)\) 实现，因为 \(PTP^{-1}=-T\)。这个式子直接显示了域自同构如何在更大的矩阵代数内成为共轭。

\begin{corollary}\label{b2:cor:csa-field-normalizer}
设 \(k\) 为域，\(A\) 为中心单 \(k\) 代数，\(L\subseteq A\) 为含 \(k\) 的子域。令
\[
N_{A^\times}(L)=\{\,u\in A^\times\mid uLu^{-1}=L\,\}.
\]
共轭给出正合列
\[
1\longrightarrow C_A(L)^\times\longrightarrow N_{A^\times}(L)
\longrightarrow\operatorname{Aut}_k(L)\longrightarrow1.
\]
若 \([L:k]=\deg A\)，则核为 \(L^\times\)。
\end{corollary}
\begin{proof}
正规化条件保证共轭限制为 \(L\) 的 \(k\) 自同构。其核恰为在 \(L\) 上逐点恒等的那些可逆元，也就是 \(C_A(L)^\times\)。对任意 \(\sigma\in\operatorname{Aut}_k(L)\)，把包含 \(\iota:L\hookrightarrow A\) 与 \(\iota\circ\sigma\) 作为两种简单代数嵌入，\cref{b2:thm:skolem-noether}给出实现它的可逆元，因此该同态满射。最后的核识别由\cref{b2:prop:csa-self-centralizing-field}。
\end{proof}

满射保证对每个 \(\sigma\in\operatorname{Aut}_k(L)\)，存在 \(u_\sigma\in N_{A^\times}(L)\) 实现它，却不保证能同时选择这些元素，使 \(u_{\sigma\tau}=u_\sigma u_\tau\) 对全部 \(\sigma,\tau\) 成立。两边虽然实现同一个域自同构，比值仍可能是中心化子中的非平凡单位。\cref{b2:prop:quaternion-normalizer-nonsplit}将在实四元数中给出没有群截面的具体正合列；下一章的循环代数还将记录实现元相乘时的这种差异。

\subsection{适用条件与非单代数反例}
\label{b2:subsec:1703:03}

先看简单子代数假设。取 \(A=M_3(k)\)、\(B=k\times k\)，定义
\[
\varphi(x,y)=\operatorname{diag}(x,y,y),\qquad
\psi(x,y)=\operatorname{diag}(x,x,y).
\]
两者都是含幺的单射代数同态，但 \(\varphi(1,0)\) 的秩为一，\(\psi(1,0)\) 的秩为二。矩阵内共轭保持矩阵的秩，所以这两个嵌入不可能由内共轭相连。此例的环境 \(A\) 中心单，失败来自 \(B\) 有两个非零简单因子。对 \(\varphi\) 的表示，两种简单分量的重数为 \(1,2\)；对 \(\psi\) 则为 \(2,1\)。总维数同为三，分量重数却不同，正是前述证明只对单 Artin 环成立的原因。

\ProseParagraphBreak
中心性也不能任意删去。把 \(\mathbb C\) 看成实代数，它是单代数，但中心大于 \(\mathbb R\)。复共轭是非恒等实代数自同构，而交换代数中的每个内自同构都是恒等映射，所以复共轭并非内自同构。

\ProseParagraphBreak
非单代数甚至可以通过交换其不同因子产生外自同构。例如 \(k\times k\) 的映射 \((x,y)\mapsto(y,x)\) 保持代数结构，却不是内自同构，因为整个代数交换。因而“自同构都是内的”不是任意有限维代数的性质；它来自中心单性及前述简单模比较的共同作用。
